\documentclass[11pt]{article}
\usepackage{fontspec}
\usepackage[french]{babel}
\usepackage{amsmath}
\usepackage{amssymb}
\usepackage{amsthm}
\usepackage{mathrsfs}
\usepackage{tikz-cd}
\usepackage{enumitem}

\newcommand{\printedpage}[1]{\marginpar{\footnotesize p.~#1}}
\newcommand{\msp}[1]{\marginpar{\footnotesize [#1]}}

% short macros for printed notation
\newcommand{\M}{\mathcal{M}}
\newcommand{\Filt}{\mathrm{Filt}}
\newcommand{\Fl}{\mathrm{Fl}}
\newcommand{\uFl}{\underline{\mathrm{Fl}}}
\newcommand{\uFilt}{\underline{\mathrm{Filt}}}
\newcommand{\card}{\operatorname{card}}
\newcommand{\Hom}{\operatorname{Hom}}
\newcommand{\Ob}{\operatorname{Ob}}
% the print has a wavy line over the expression (isomorphism classes);
% rendered here by a wide tilde
\newcommand{\wavy}[1]{\widetilde{#1}}

\begin{document}

\printedpage{14}
% head: A. Grothendieck : Dérivateurs, Ch. XV, Édition des Universités de Montpellier II et Paris VII
% foot: Version 5 novembre 2010
\msp{page 20}
première formule (3.7 bis), notons que l'on a égalité lorsque $\M$ est la catégorie des couples $(A,B)$, $A$ un ensemble et $B$ une partie de $A$, avec les flèches évidentes, et qu'on définit $\M^\pi$ (pour $\pi\geq\pi_0$) comme la sous-catégorie strictement pleine formée des $(A,B)$ tels que $\card A\geq\pi$.

\textsc{Démonstration de} 3.1. Prouvons d'abord la formule (3.8),
\[
\card\Hom(A,X)\ \leq\ s^{(r^{\pi_0})}\qquad\quad
\left\{\begin{array}{l}
\text{si } A\in\Filt^r,\ X\in\Filt^s,\\[1.5ex]
r\geq\pi_1,\ s\geq\sup(\pi_1,C_1).
\end{array}\right.
\]

(\textbf{N.B.} Je n'utiliserai pas $r\geq C_1$, postulé dans 3.1 $1^\circ$.) J'écris \texttt{[<<J'écris>> rayé]}
\[
\begin{aligned}
A\ &=\ \varinjlim_I A_i\ ,\quad \card I\ \leq\ r^{\pi_0}\ ,\quad A_i\ \in\ \Filt^{\pi_0}\ ,\\
X\ &=\ \varinjlim_J X_j\ ,\quad \card J\ \leq\ s^{\pi_1}\ ,\quad X_j\ \in\ \Filt^{\pi_1}\ ,\qquad\qquad J\text{ grand devant }\pi_1\ .
\end{aligned}
\]
On aura
\[
\Hom(A,X)\ \simeq\ \varprojlim_{i\in I}\ \underbrace{\Hom(A_i,\,\varinjlim_J X_j)}_{\overset{(\text{\footnotemark[6]})}{\simeq}\ \varinjlim\limits_{j\in J}\ \underbrace{\Hom(A_i,X_j)}_{\text{de cardinal }\leq C_1}}\ ,
\]
\footnotetext[6]{car les $A\in\Filt^{\pi_0}$ sont $\pi_1$-accessibles.}
donc pour tout $i\in I$
\[
\Hom(A_i,X)\ \leq\ \card(J)\times C_1\ \leq\ s^{\pi_1}\times C_1\ =\ s^{\pi_1}\ ,
\]
car $s\geq C_1$, donc $s^{\pi_1}\geq C_1$, d'où
\[
\card(\,\varprojlim_{j\in I}\Hom(A_i,X))\ \leq\ (s^{\pi_1})^{\card I}\ \leq\ \underbrace{(s^{\pi_1})^{(r^{\pi_0})}}_{s^{\pi_1 r^{\pi_0}}}\ . % [sic] j\in I -> i\in I, the limit is over the index i of A_i
\]
Comme $r\geq\pi_1$, on a $r^{\pi_0}\geq\pi_1$, donc $r^{\pi_0}\pi_1=r^{\pi_0}$, d'où
\[
\card\Hom(A,X)\ \leq\ s^{(r^{\pi_0})}\ ,
\]
c'est bien la formule (3.8). Pour en tirer une estimation pour $\M^r/X$ (avec $X\in\Filt^s$), \msp{page 21} je vais d'abord majorer le cardinal des classes d'isomorphie d'objets de $\M^r/X$, \textit{i.e.} des flèches à isomorphisme près (dans $\M^r$)
\[
A\ \longrightarrow\ X\ ,\qquad\quad X\text{ fixé dans }\Filt^s,\ A\text{ variable dans }\M^r\ .
\]
On aura
\[
A\ =\ \varinjlim_{i\in I} A_i\qquad\quad\text{avec } A_i\in\Filt^{\pi_0},\ I\text{ ensemble ordonné de cardinal}\leq r^{\pi_0}\ .
\]

\printedpage{15}
% head: A. Grothendieck : Dérivateurs, Ch. XV, Édition des Universités de Montpellier II et Paris VII
% foot: Version 5 novembre 2010
Pour $I$ fixé, le cardinal de l'ensemble des systèmes inductifs indexés par $I$, dans $\M^{\pi_0}/X$, à isomorphisme près, peut se majorer ainsi : pour tout $i\in I$, il faut choisir $A_i$ dans $\M^{\pi_0}$ (mais seule sa classe dans $\wavy{\M^{\pi_0}}$ importe), donc c'est majoré par card $\wavy{\M^{\pi_0}}\ \leq C_0$~\textsuperscript{(}\footnote[7]{On a posé ici $C_0=\card\wavy{\uFl^{\pi_0}}$ , donc $C_0\leq C_1$.}\textsuperscript{)}, et pour tous les $i$, par $\card I\times C_0\leq r^{\pi_0}\times C_0$, et si $r\geq C_0$, donc $r^{\pi_0}\geq C_0$, on trouve une majoration des choix possibles du système des \textit{objets} $A_i$ dans $\M^\pi$, à isomorphisme près, par $r^{\pi_0}$. Si on se donne de plus les morphismes de transition $A_i\longrightarrow A_j$ pour $i<j$, cela fait pour chaque paire comparable, élément de $I\times I$, un choix majoré par $C_0$, donc encore une majoration par $\ \underbrace{\card(I\times I)}_{=(\card I)^2=\card I}\cdot C_0\leq r^{\pi_0}\cdot C_0=r^{\pi_0}$. Donc l'ensemble des systèmes inductifs en question, à isomorphisme près, pour $I$ fixé, est de cardinal $\leq r^{\pi_0}\cdot r^{\pi_0}=r^{\pi_0}$. Enfin, les \textit{homomorphismes} d'un système inductif donné dans $X$ sont majorés par $\card I\times s^{(r^{\pi_0})}$ (par (3.8) déjà prouvé), donc par
\[
r^{\pi_0}\cdot s^{(r^{\pi_0})}\ =\ s^{(r^{\pi_0})}\qquad\quad(\text{puisque } r^{\pi_0}\leq s^{(r^{\pi_0})})\ .
\]

Au total, pour $I$ fixé, l'ensemble des systèmes inductifs à isomorphisme près, indexés par $I$, à valeurs dans $\M^r/X$, est \msp{page 22} de cardinal majoré par $\quad\underbrace{r^{\pi_0}}_{\text{\scriptsize\begin{tabular}{@{}c@{}}nombre des\\ systèmes inductifs\\ dans $\M^\pi$, à\\ isomorphisme près\end{tabular}}}\quad\cdot\quad\underbrace{s^{(r^{\pi_0})}}_{\text{\scriptsize\begin{tabular}{@{}c@{}}nombre des\\ flèches d'un\\ tel système inductif\\ dans $X$\end{tabular}}}\quad=s^{(r^{\pi_0})}$.

Mais il faut maintenant faire varier $I$, à isomorphisme près, dans la catégorie des ensembles ordonnés (filtrants \textit{etc.}) de cardinal $\leq\lambda=r^{\pi_0}$. On voit de suite que cet ensemble de classes d'isomorphie (qui peut s'interpréter comme un ensemble quotient d'un certain ensemble de parties d'un ensemble fixe $I_0$ de cardinal $\lambda$) est de cardinal $\leq 2^\lambda$. On trouve donc finalement
\[
\card\wavy{\M^r/X}\ \ \leq\ s^{(r^{\pi_0})}\cdot 2^{(r^{\pi_0})}\ =\ s^{(r^{\pi_0})}\ ,
\]
sous condition que
\[
r\ \geq\ \sup(\pi_1,C_0)\ ,\qquad\quad s\ \geq\ \sup(\pi_1,C_1)\ .
\]
Il faut prouver l'inégalité similaire pour card $\wavy{\uFl^r\,\M}$ \ (3.9). Or, comme pour $A$, $B$ dans $\M^r$, on a, si $r\geq\sup(\pi_1,C_1)$,
\[
\card\Hom_{\M}(A,B)\ \leq\ r^{(r^{\pi_0})}
\]
en vertu de (3.8) où on fait $s=r$, on trouve que
\[
\begin{aligned}
\card\wavy{\uFl(\M^r/X)}\ &\leq\ \Bigl(\card\wavy{\M^r/X}\,\Bigr)^2\cdot r^{(r^{\pi_0})}\\
&\leq\ s^{(r^{\pi_0})}\cdot r^{(r^{\pi_0})}\\
&=\ s^{(r^{\pi_0})}\qquad\text{si } s\geq r\ .
\end{aligned}
\]
(On va regarder plus bas les cas $s\leq r$~\textsuperscript{(}\footnote[8]{\textit{cf.} cor. 2, page 23.}\textsuperscript{)}). C'est bien la formule (3.9).

\printedpage{16}
% head: A. Grothendieck : Dérivateurs, Ch. XV, Édition des Universités de Montpellier II et Paris VII
% foot: Version 5 novembre 2010
\msp{page 23}
Venons-en à la catégorie $\M^r$ elle-même. L'argument de la page 21 montre que l'ensemble des classes d'isomorphie d'objets de $\M^r$ indexés par un $I$ fixé, de cardinal $\leq r^{\pi_0}$, comme $A=\varinjlim\limits_i A_i$ avec $A_i\in\Filt^{\pi_0}$, est de cardinal majoré par $r^{\pi_0}$ (si $r\geq C_0$). Si on fait varier $I$, on trouve une majoration par $r^{\pi_0}\cdot 2^{(r^{\pi_0})}=2^{(r^{\pi_0})}$. On en conclut comme tantôt que l'ensemble $\ \wavy{\Fl\,\M^r}$ est de cardinal majoré par
\[
(2^{(r^{\pi_0})})^2\times r^{(r^{\pi_0})}\qquad\quad(\text{si } r\geq\sup(\pi_1,C_1))\ ,
\]
donc par $r^{(r^{\pi_0})}$. Je dis qu'on a
\[
r^{(r^{\pi_0})}\ =\ 2^{(r^{\pi_0})}\ ,
\]
en effet, on a $\geq$, et pour prouver $\leq$, on note que $r<2^r$, donc $r^{(r^{\pi_0})}\leq(2^r)^{(r^{\pi_0})}=2^{r\cdot r^{\pi_0}}=2^{(r^{\pi_0})}$, O.K.! Cela établit donc 3.7. Cet argument prouve en fait le

\textbf{Lemme 3.2.} \textit{Soient $r$ un cardinal, $r'$ un cardinal $\geq r$. Alors on a}
\[
r^{r'}\ =\ 2^{r'}\qquad\quad(=s^{r'}\ \textit{pour tout}\ 2\leq s\leq r)\ .
\]

\textbf{Corollaire 1.} \textit{Soient $r\geq\sup(\pi_1,C_1)$, et $A,X\in\Filt^\pi$. Alors}
\[
\card\Hom_{\M}(A,X)\ \leq\ 2^{(r^{\pi_0})}\qquad\quad(=2^r\ \textit{si}\ r=2^c,\ \textit{avec}\ c\geq\pi_0)\ . \tag*{(3.10)}
\]

On applique (3.8), qui donne la majoration $r^{(r^{\pi_0})}$, et le lemme 3.2.

On notera qu'on trouve la même majoration pour le cardinal de $\wavy{\uFl(\M^r)}$ tout entier (3.7)! De même, (3.9) donne :

\textbf{Corollaire 2.} \textit{Soit $r\geq\sup(\pi_1,C_1)$, et $X\in\Ob\M^r$. Alors}
\[
\card\wavy{\uFl(\M^r/X)}\ \leq\ 2^{(r^{\pi_0})}\qquad\quad(=2^r\ \textit{si}\ r=2^c,\ \textit{avec}\ c\geq\pi_0)\ . \tag*{(3.11)}
\]
c'est donc la formule (3.9) pour $s\leq r$, laissée en suspens.

\msp{page 24}
Revenons à la majoration cardinale des $X_\alpha(f)$, où $f : X\longrightarrow Y$ est dans $\Fl^s(\M)$, \textit{i.e.} $X,Y\in\Filt^s$. On suppose de plus $C_0\subset\Fl^r(\M)$,
\[
\begin{aligned}
&C_0\ \subset\ \Fl^r(\M)\ =\ \Fl(\M^r)\ ,\qquad\quad(f:X\longrightarrow Y)\in\Fl^s(\M)\ ,\ \ \textit{i.e.}\ X,Y\in\Filt^s\ ,\\
&\text{avec } s\geq r\ .
\end{aligned} \tag*{(3.10)}
\]

\texttt{[Deux fois (3.10).]} Il y a une difficulté idiote, c'est que $\M^r$ est équivalente à une petite catégorie, mais n'est pas elle-même petite. Même si $C_0$ ne contenait qu'une seule classe d'isomorphie de flèches de $\M^r$, son cardinal pourrait être (pour $r$ fixé) arbitrairement élevé. Aussi, dans ce qui suit, je dois supposer la condition de “non pléthoricité cardinale” pour $C_0$ :
\[
\parbox{0.72\linewidth}{$\forall\,i\in C_0$, désignant par $C_0(i)$ l'ensemble des $i'\in C_0$ telles que $i\simeq i'$ dans $\Fl(\M)$, j'impose que $\card C_0(i)\leq 2^{(r^{\pi_0})}$ (ce qui est aussi la valeur de la majoration pour le cardinal de $\wavy{\uFl^r(\M)}=\wavy{\uFl(\M^r)}$\,).} \tag*{(3.11)}
\]

\printedpage{17}
% head: A. Grothendieck : Dérivateurs, Ch. XV, Édition des Universités de Montpellier II et Paris VII
% foot: Version 5 novembre 2010
\texttt{[Deux fois (3.11).]} Cela implique donc que l'on a $\card C_0\leq(2^{(r^{\pi_0})})^2=2^{(r^{\pi_0})}$,
\[
\card C_0\ \leq\ 2^{(r^{\pi_0})}\ , \tag*{(3.12)}
\]
qui est équivalente à (3.11) (moyennant $C_0\subset\Fl^r(\M)$, postulé dans (3.10)).

Il s'ensuit par (3.8) que l'on a
\[
\card\Delta(C_0,f)\ \leq\ \underbrace{\card(C_0)}_{\leq 2^{(r^{\pi_0})}}\cdot s^{(r^{\pi_0})}\cdot s^{(r^{\pi_0})}\ \leq\ s^{(r^{\pi_0})}\ .
\]

D'autre part, les $X_D$ dont on prend la somme amalgamée sous $X$, pour obtenir $X_1(f)$, sont dans $\M^s$, si on fait l'hypothèse suivante. \msp{page 25}
\[
\left|\ \parbox{0.72\linewidth}{Les $\M^\pi=\uFilt^\pi$ sont stables par sommes amalgamées simples, et par $\varinjlim$ filtrantes de cardinal $\leq\pi$ (pas nécessairement grandes devant $\pi_0$) ,}\right. \tag*{(3.13)}
\]
compte tenu que $r\leq s$, donc $\M^r\subset\M^s$, donc $A$, $B$, $X$ dans $X_D=B\amalg_A X$ sont dans $\M^s$. Les sommes amalgamées partielles \textit{finies} des $X_D$ sont donc aussi dans $\M^s$, par récurrence sur le cardinal de l'ensemble d'indices. La somme amalgamée infinie apparaît ainsi comme $\varinjlim$ filtrante d'objets de $\M^s$, l'ensemble d'indices étant l'ensemble $\mathfrak{P}^*_{\mathfrak{f}}\Delta(C_0,f)$ des parties finies non vides de $\Delta(C_0,f)$. Cet ensemble est de cardinal majoré par $s'$, si $s'$ est un cardinal \textit{infini} tel que $\card\Delta(C_0,f)\leq s'$. On peut donc prendre $s'=s^{(r^{\pi_0})}$. On a donc prouvé, compte tenu de (3.13) appliqué à $\M^{s'}$, le % ?? uncertain: symbol in X_D = B \amalg_A X -- \amalg or a doubled \sqcup; subscript of \mathfrak{P}^* -- fraktur f or roman f

\textbf{Lemme 3.3.} \textit{On suppose} \textup{(3.10), i.e.} $C_0\subset\Fl^r(\M)$, $f\in\Fl^s(\M)$, \textit{avec} $s\geq r\geq\sup(\pi_1,C_1)$, \textit{et de plus que $C_0$ satisfait la condition de non-pléthoricité} \textup{(3.11)} \textup{(}\textit{ce qu'on peut toujours supposer, quitte à réduire $C_0$ sans changer $\widetilde{C_0}=\Phi$ et $C_{0*}=\Psi$}\textup{)}\textit{, et que la filtration cardinale satisfait la condition de stabilité (anodine!)} \textup{(3.13)}\textit{. Posons}
\[
s'\ =\ s^{(r^{\pi_0})}\ . \tag*{(3.14)}
\]
\textit{Alors $X_1(C_0,f)$ est dans $\M^{s'}$. Plus généralement et plus joliment, pour $C_0$ et $r$ fixés} \textup{(}\textit{avec $C_0\subset\Fl^r(\M)$ et satisfaisant la condition de non-pléthoricité} \textup{(3.11)} \textit{relativement à $r$}\textup{)}\textit{, pour tout cardinal} \msp{page 26} \textit{de la forme}
\[
s'\ =\ s^t\ ,\qquad\quad\textit{avec}\ t\geq r^{\pi_0},\ s\geq 2 \tag*{(3.15)}
\]
\textit{(d'où $s'\geq 2^{(r^{\pi_0})}>r^{\pi_0}\geq r$), $\uFl^{s'}(\M)$ est stable sous le foncteur $p_1^{C_0}:\uFl(\M)\longrightarrow\uFl(\M)$.}

En effet, si on applique la première partie du lemme au couple $(r,s')$, on trouve que $X_1(f)$ est dans $\M^{s''}$, avec $s''=s'^{(r^{\pi_0})}$. Mais
\[
s''\ =\ s'^{(r^{\pi_0})}\ =\ (s^t)^{(r^{\pi_0})}\ =\ s^{tr^{\pi_0}}\ \underset{\substack{\vert\\ \text{\scriptsize car $t\geq r^{\pi_0}$,}\\ \text{\scriptsize donc $tr^{\pi_0}=t$}}}{=}\ s^t\ =\ s'\ ,
\]
\printedpage{18}
% head: A. Grothendieck : Dérivateurs, Ch. XV, Édition des Universités de Montpellier II et Paris VII
% foot: Version 5 novembre 2010
donc on voit que $f\in\Fl^{s'}$ implique $X_1(f)\in\M^{s'}$, ou ce qui revient au même, $p_1(f)\in\Fl^{s'}$.

\textbf{Corollaire.} \textit{Soient $C_0$, $r$, $s'=s^t$ $(t\geq r^{\pi_0},\ s\geq 2)$ comme dessus, et soit $\alpha$ un nombre ordinal tel que}
\[
\card\alpha\ \leq\ s'\ =\ s^t\ .
\]
\textit{Alors $\uFl^{s'}(\M)$ est stable par le foncteur $p_\alpha=p_\alpha^{C_0}$.}

On le voit aussitôt par récurrence transfinie, compte tenu du lemme 3.3, et de l'hypothèse de stabilité (3.13).

Pour résumer :

\textbf{Théorème 2.} \textit{Soit $\M$ une catégorie stable par petites $\varinjlim$ filtrantes, et munie d'une filtration cardinale} $(\Filt^\pi)_{\pi\geq\pi_0}$\textit{, telle que les} $\Filt^\pi$ \textit{soient stables par $\varinjlim$ filtrantes de cardinal $\leq\pi$ (pas forcément grandes devant $\pi_0$), et par les sommes amalgamées simples qui existent dans $\M$. On suppose que les éléments de} $\Filt^{\pi_0}$ \textit{soient accessibles, et on désigne par $\pi_1$ le plus petit cardinal $\geq\pi_0$, tel que les objets de} $\Filt^{\pi_0}$ \textit{soient $\pi_1$-accessibles. On pose} \msp{page 27}
\[
c_1\ =\ \card\Fl(\M^{\pi_1})\ =\ \card\Fl^{\pi_1}(\M)
\]
\textit{(où pour tout $\pi\geq\pi_0$, $\M^\pi$ désigne la sous-catégorie strictement pleine de $\M$ définie par} $\Filt^\pi$\textit{, et} $\Fl^\pi(\M)=\Fl(\M^\pi)$\textit{).}

\textit{Soit $C_0$ une petite partie de} $\Fl(\M)$\textit{, formée de flèches coquarrables, et posons}
\[
\Phi\ =\ \widetilde{C_0}\ =\ \Psi^*\ ,\qquad\quad\textit{où}\ \ \Psi=C_{0*}=\Phi_*\ .
\]
\textit{Soit $r$ un cardinal} $\geq\sup(\pi_1,c_1)$\textit{, et tel que l'on ait $C_0\subset\Fl^r(\M)$. On suppose que $C_0$ est irrédondant (}i.e.\textit{ deux objets de $C_0$ isomorphes sont égaux), ou plus généralement, que $C_0$ satisfait la très anodine condition de non-pléthoricité} \textup{(3.11)} \textit{relativement à $r$. On considère la suite transfinie des foncteurs $p_\alpha=p_\alpha^{C_0}$ sur $\uFl(\M)$} \textup{(}\textit{cf.} \textup{th. 1)}\textit{, donnant lieu aux factorisations fonctorielles d'une $f\in\uFl(\M)$, $f:X\longrightarrow Y$, en}
\[
X\ \xrightarrow{i_\alpha(f)}\ X_\alpha(f)\ \xrightarrow{p_\alpha(f)}\ Y\ . \tag*{(3.16)}
\]

\textit{On a alors ce qui suit :}

$1^\circ$) \textit{On a (pour mémoire) $i_\alpha(f)\in\Phi$ pour tout $\alpha$, et de plus,} si $\alpha$ est grand devant $r^{\pi_0}$, \textit{on a aussi $p_\alpha(f)\in\Psi$. Si $r'$ est un cardinal $>r^{\pi_0}$ (p. ex. le cardinal successeur de $r^{\pi_0}$), on peut prendre pour $\alpha$ le premier ordinal dont le cardinal soit égal à $r'$.}

$2^\circ$) \textit{Pour tout cardinal $s$ de la forme}
\[
s\ =\ c^t\qquad\quad\textit{avec}\ \ t\geq r^{\pi_0},\ c\geq 2
\]
\textit{(donc $s=2^t$ si on prend $c\leq t$,} cf. \textit{lemme 3.2), $\uFl^s(\M)$ est stable sous les $p_\alpha$ tels que} $\card\alpha\leq s$\textit{. Si on prend p. ex. pour $\alpha$ le premier ordinal dont le cardinal soit $>r^{\pi_0}$, alors $\uFl^s(\M)$ est stable par $p_\alpha$, qui par $1^\circ$) réalise donc une factorisation fonctorielle} \textup{(3.16)} \textit{avec $i_\alpha(f)\in\Phi$, $p_\alpha(f)\in\Psi$.}

\end{document}
